
\documentclass{article}
\usepackage{colortbl}
\usepackage{makecell}
\usepackage{multirow}
\usepackage{supertabular}

\begin{document}

\newcounter{utterance}

\centering \large Interaction Transcript for game `matchit\_ascii', experiment `similar\_grid\_1', episode 1 with Qwen3.6{-}35B{-}A3B{-}FP8{-}without{-}reasoning.
\vspace{24pt}

{ \footnotesize  \setcounter{utterance}{1}
\setlength{\tabcolsep}{0pt}
\begin{supertabular}{c@{$\;$}|p{.15\linewidth}@{}p{.15\linewidth}p{.15\linewidth}p{.15\linewidth}p{.15\linewidth}p{.15\linewidth}}
    \# & $\;$A & \multicolumn{4}{c}{Game Master} & $\;\:$B\\
    \hline

    \theutterance \stepcounter{utterance}  
    & & \multicolumn{4}{p{0.6\linewidth}}{
        \cellcolor[rgb]{0.9,0.9,0.9}{
            \makecell[{{p{\linewidth}}}]{
                \texttt{\tiny{[P1$\langle$GM]}}
                \texttt{Участвате в съвместна игра на отгатване.} \\
\texttt{Целта е да разберете дали тази мрежа и друга картина, която виждам само аз, са еднакви.} \\
\texttt{Моля, първо опишете вашата мрежа. След това аз ще дам моето описание и ще можем да си задаваме взаимно въпроси за мрежите, за да разберем дали са еднакви. Сега започнете краткото си описание на мрежата с "ОПИСАНИЕ:", последвано от описанието. Не добавяйте нищо друго. Това е вашата мрежа:} \\
\texttt{X X X X X} \\
\texttt{X X X X ▢} \\
\texttt{X X X ▢ ▢} \\
\texttt{X X ▢ ▢ ▢} \\
\texttt{X ▢ ▢ ▢ ▢} \\
            }
        }
    }
    & \\ \\

    \theutterance \stepcounter{utterance}  
    & \multicolumn{4}{p{0.6\linewidth}}{
        \cellcolor[rgb]{0.9,0.9,0.9}{
            \makecell[{{p{\linewidth}}}]{
                \texttt{\tiny{[P1$\rangle$GM]}}
                \texttt{ОПИСАНИЕ: Мрежата е квадратна с размери 5x5. Състои се от символи "X" и празни квадрати "▢". Символите "X" запълват горния ляв триъгълник, включително главния диагонал, докато празните квадрати "▢" запълват долния десен триъгълник, без да включват главния диагонал. По{-}конкретно:} \\
\texttt{{-} Ред 1: X X X X X} \\
\texttt{{-} Ред 2: X X X X ▢} \\
\texttt{{-} Ред 3: X X X ▢ ▢} \\
\texttt{{-} Ред 4: X X ▢ ▢ ▢} \\
\texttt{{-} Ред 5: X ▢ ▢ ▢ ▢} \\
            }
        }
    }
    & & \\ \\

    \theutterance \stepcounter{utterance}  
    & & & \multicolumn{2}{p{0.3\linewidth}}{
        \cellcolor[rgb]{0.9,0.9,0.9}{
            \makecell[{{p{\linewidth}}}]{
                \texttt{\tiny{[GM$|$GM]}}
                \texttt{continue} \\
            }
        }
    }
    & & \\ \\

    \theutterance \stepcounter{utterance}  
    & & \multicolumn{4}{p{0.6\linewidth}}{
        \cellcolor[rgb]{0.9,0.9,0.9}{
            \makecell[{{p{\linewidth}}}]{
                \texttt{\tiny{[GM$\rangle$P2]}}
                \texttt{Участвате в съвместна игра на отгатване.} \\
\texttt{Целта е да разберете дали тази мрежа и друга картина, която виждам само аз, са еднакви.} \\
\texttt{Моля, първо опишете вашата мрежа. След това аз ще дам моето описание и ще можем да си задаваме взаимно въпроси за мрежите, за да разберем дали са еднакви. Сега започнете краткото си описание на мрежата с "ОПИСАНИЕ:", последвано от описанието. Не добавяйте нищо друго. Това е вашата мрежа:} \\
\texttt{X X X X X} \\
\texttt{▢ X X X X} \\
\texttt{▢ ▢ X X X} \\
\texttt{▢ ▢ ▢ X X} \\
\texttt{▢ ▢ ▢ ▢ X} \\
            }
        }
    }
    & \\ \\

    \theutterance \stepcounter{utterance}  
    & & & \multicolumn{4}{p{0.6\linewidth}}{
        \cellcolor[rgb]{0.9,0.9,0.9}{
            \makecell[{{p{\linewidth}}}]{
                \texttt{\tiny{[GM$\langle$P2]}}
                \texttt{ОПИСАНИЕ: Мрежата е с размери 5x5. В горния ляв ъгъл има триангуларна форма от празни клетки (▢), които запълват долната лява част на матрицата. Конкретно, първият ред е напълно запълнен с X{-}ове. Вторият ред има една празна клетка в началото, следвана от четири X{-}а. Третият ред има две празни клетки в началото, следвани от три X{-}а. Четвъртият ред има три празни клетки в началото, следвани от два X{-}а. Петият ред има четири празни клетки в началото, следвани от един X. Общо има 15 клетки с X и 10 празни клетки (▢).} \\
            }
        }
    }
     \\ \\

    \theutterance \stepcounter{utterance}  
    & & & \multicolumn{2}{p{0.3\linewidth}}{
        \cellcolor[rgb]{0.9,0.9,0.9}{
            \makecell[{{p{\linewidth}}}]{
                \texttt{\tiny{[GM$|$GM]}}
                \texttt{continue} \\
            }
        }
    }
    & & \\ \\

    \theutterance \stepcounter{utterance}  
    & & \multicolumn{4}{p{0.6\linewidth}}{
        \cellcolor[rgb]{0.9,0.9,0.9}{
            \makecell[{{p{\linewidth}}}]{
                \texttt{\tiny{[GM$\rangle$P2]}}
                \texttt{Това е моето ОПИСАНИЕОПИСАНИЕ: Мрежата е квадратна с размери 5x5. Състои се от символи "X" и празни квадрати "▢". Символите "X" запълват горния ляв триъгълник, включително главния диагонал, докато празните квадрати "▢" запълват долния десен триъгълник, без да включват главния диагонал. По{-}конкретно:} \\
\texttt{{-} Ред 1: X X X X X} \\
\texttt{{-} Ред 2: X X X X ▢} \\
\texttt{{-} Ред 3: X X X ▢ ▢} \\
\texttt{{-} Ред 4: X X ▢ ▢ ▢} \\
\texttt{{-} Ред 5: X ▢ ▢ ▢ ▢} \\
\texttt{Сега задайте въпрос, за да разберете нови аспекти на моята мрежа, които може да се различават от вашата мрежа. Започнете с "ВЪПРОС:" и не добавяйте нищо друго.} \\
            }
        }
    }
    & \\ \\

    \theutterance \stepcounter{utterance}  
    & & & \multicolumn{4}{p{0.6\linewidth}}{
        \cellcolor[rgb]{0.9,0.9,0.9}{
            \makecell[{{p{\linewidth}}}]{
                \texttt{\tiny{[GM$\langle$P2]}}
                \texttt{ВЪПРОС: В моята мрежа символите "X" запълват горния десен триъгълник (включително главния диагонал), а празните клетки "▢" са в долния ляв триъгълник. Вашата мрежа ли е с "X" в горния ляв триъгълник и "▢" в долния десен?} \\
            }
        }
    }
     \\ \\

    \theutterance \stepcounter{utterance}  
    & & & \multicolumn{2}{p{0.3\linewidth}}{
        \cellcolor[rgb]{0.9,0.9,0.9}{
            \makecell[{{p{\linewidth}}}]{
                \texttt{\tiny{[GM$|$GM]}}
                \texttt{continue} \\
            }
        }
    }
    & & \\ \\

    \theutterance \stepcounter{utterance}  
    & & \multicolumn{4}{p{0.6\linewidth}}{
        \cellcolor[rgb]{0.9,0.9,0.9}{
            \makecell[{{p{\linewidth}}}]{
                \texttt{\tiny{[P1$\langle$GM]}}
                \texttt{Това е моето ОПИСАНИЕОПИСАНИЕ: Мрежата е с размери 5x5. В горния ляв ъгъл има триангуларна форма от празни клетки (▢), които запълват долната лява част на матрицата. Конкретно, първият ред е напълно запълнен с X{-}ове. Вторият ред има една празна клетка в началото, следвана от четири X{-}а. Третият ред има две празни клетки в началото, следвани от три X{-}а. Четвъртият ред има три празни клетки в началото, следвани от два X{-}а. Петият ред има четири празни клетки в началото, следвани от един X. Общо има 15 клетки с X и 10 празни клетки (▢).} \\
\texttt{ВЪПРОС: В моята мрежа символите "X" запълват горния десен триъгълник (включително главния диагонал), а празните клетки "▢" са в долния ляв триъгълник. Вашата мрежа ли е с "X" в горния ляв триъгълник и "▢" в долния десен?Започнете отговора си с "ОТГОВОР:" и не добавяйте нищо друго.} \\
            }
        }
    }
    & \\ \\

    \theutterance \stepcounter{utterance}  
    & \multicolumn{4}{p{0.6\linewidth}}{
        \cellcolor[rgb]{0.9,0.9,0.9}{
            \makecell[{{p{\linewidth}}}]{
                \texttt{\tiny{[P1$\rangle$GM]}}
                \texttt{ОТГОВОР: Не, нашите мрежи не са еднакви. В моята мрежа символите "X" са в горния ляв триъгълник (включително главния диагонал), а празните клетки "▢" са в долния десен триъгълник. В вашата мрежа, според описанието, "X" са в горния десен триъгълник, а "▢" са в долния ляв. Това означава, че позициите на символите са огледално отражение едно на друго спряо главния диагонал.} \\
            }
        }
    }
    & & \\ \\

    \theutterance \stepcounter{utterance}  
    & & & \multicolumn{2}{p{0.3\linewidth}}{
        \cellcolor[rgb]{0.9,0.9,0.9}{
            \makecell[{{p{\linewidth}}}]{
                \texttt{\tiny{[GM$|$GM]}}
                \texttt{continue} \\
            }
        }
    }
    & & \\ \\

    \theutterance \stepcounter{utterance}  
    & & \multicolumn{4}{p{0.6\linewidth}}{
        \cellcolor[rgb]{0.9,0.9,0.9}{
            \makecell[{{p{\linewidth}}}]{
                \texttt{\tiny{[P1$\langle$GM]}}
                \texttt{Сега задайте въпрос, за да разберете нови аспекти на моята мрежа, които може да се различават от вашата мрежа. Започнете с "ВЪПРОС:" и не добавяйте нищо друго.} \\
            }
        }
    }
    & \\ \\

    \theutterance \stepcounter{utterance}  
    & \multicolumn{4}{p{0.6\linewidth}}{
        \cellcolor[rgb]{0.9,0.9,0.9}{
            \makecell[{{p{\linewidth}}}]{
                \texttt{\tiny{[P1$\rangle$GM]}}
                \texttt{ВЪПРОС: Ако транспонирам моята мрежа (разменям редовете със стълбовете), дали ще получа точно вашата мрежа?} \\
            }
        }
    }
    & & \\ \\

    \theutterance \stepcounter{utterance}  
    & & & \multicolumn{2}{p{0.3\linewidth}}{
        \cellcolor[rgb]{0.9,0.9,0.9}{
            \makecell[{{p{\linewidth}}}]{
                \texttt{\tiny{[GM$|$GM]}}
                \texttt{continue} \\
            }
        }
    }
    & & \\ \\

    \theutterance \stepcounter{utterance}  
    & & \multicolumn{4}{p{0.6\linewidth}}{
        \cellcolor[rgb]{0.9,0.9,0.9}{
            \makecell[{{p{\linewidth}}}]{
                \texttt{\tiny{[GM$\rangle$P2]}}
                \texttt{ОТГОВОР: Не, нашите мрежи не са еднакви. В моята мрежа символите "X" са в горния ляв триъгълник (включително главния диагонал), а празните клетки "▢" са в долния десен триъгълник. В вашата мрежа, според описанието, "X" са в горния десен триъгълник, а "▢" са в долния ляв. Това означава, че позициите на символите са огледално отражение едно на друго спряо главния диагонал.} \\
\texttt{ВЪПРОС: Ако транспонирам моята мрежа (разменям редовете със стълбовете), дали ще получа точно вашата мрежа?Започнете отговора си с "ОТГОВОР:" и не добавяйте нищо друго.} \\
            }
        }
    }
    & \\ \\

    \theutterance \stepcounter{utterance}  
    & & & \multicolumn{4}{p{0.6\linewidth}}{
        \cellcolor[rgb]{0.9,0.9,0.9}{
            \makecell[{{p{\linewidth}}}]{
                \texttt{\tiny{[GM$\langle$P2]}}
                \texttt{ОТГОВОР: Да, ако транспонирате вашата мрежа, ще получите точно моята мрежа. Вашата мрежа има "X" в горния ляв триъгълник (включително диагонала), което при транспониране се превръща в "X" в горния десен триъгълник (включително диагонала), което съвпада с моето описание.} \\
            }
        }
    }
     \\ \\

    \theutterance \stepcounter{utterance}  
    & & & \multicolumn{2}{p{0.3\linewidth}}{
        \cellcolor[rgb]{0.9,0.9,0.9}{
            \makecell[{{p{\linewidth}}}]{
                \texttt{\tiny{[GM$|$GM]}}
                \texttt{continue} \\
            }
        }
    }
    & & \\ \\

    \theutterance \stepcounter{utterance}  
    & & \multicolumn{4}{p{0.6\linewidth}}{
        \cellcolor[rgb]{0.9,0.9,0.9}{
            \makecell[{{p{\linewidth}}}]{
                \texttt{\tiny{[GM$\rangle$P2]}}
                \texttt{Сега задайте въпрос, за да разберете нови аспекти на моята мрежа, които може да се различават от вашата мрежа. Започнете с "ВЪПРОС:" и не добавяйте нищо друго.} \\
            }
        }
    }
    & \\ \\

    \theutterance \stepcounter{utterance}  
    & & & \multicolumn{4}{p{0.6\linewidth}}{
        \cellcolor[rgb]{0.9,0.9,0.9}{
            \makecell[{{p{\linewidth}}}]{
                \texttt{\tiny{[GM$\langle$P2]}}
                \texttt{ВЪПРОС: Ако завъртите вашата мрежа на 180 градуса, дали ще получите точно моята мрежа?} \\
            }
        }
    }
     \\ \\

    \theutterance \stepcounter{utterance}  
    & & & \multicolumn{2}{p{0.3\linewidth}}{
        \cellcolor[rgb]{0.9,0.9,0.9}{
            \makecell[{{p{\linewidth}}}]{
                \texttt{\tiny{[GM$|$GM]}}
                \texttt{continue} \\
            }
        }
    }
    & & \\ \\

    \theutterance \stepcounter{utterance}  
    & & \multicolumn{4}{p{0.6\linewidth}}{
        \cellcolor[rgb]{0.9,0.9,0.9}{
            \makecell[{{p{\linewidth}}}]{
                \texttt{\tiny{[P1$\langle$GM]}}
                \texttt{ОТГОВОР: Да, ако транспонирате вашата мрежа, ще получите точно моята мрежа. Вашата мрежа има "X" в горния ляв триъгълник (включително диагонала), което при транспониране се превръща в "X" в горния десен триъгълник (включително диагонала), което съвпада с моето описание.} \\
\texttt{ВЪПРОС: Ако завъртите вашата мрежа на 180 градуса, дали ще получите точно моята мрежа?Започнете отговора си с "ОТГОВОР:" и не добавяйте нищо друго.} \\
            }
        }
    }
    & \\ \\

    \theutterance \stepcounter{utterance}  
    & \multicolumn{4}{p{0.6\linewidth}}{
        \cellcolor[rgb]{0.9,0.9,0.9}{
            \makecell[{{p{\linewidth}}}]{
                \texttt{\tiny{[P1$\rangle$GM]}}
                \texttt{ОТГОВОР: Не, ако завъртите моята мрежа на 180 градуса, няма да получите вашата мрежа. Завъртането на 180 градуса на моята мрежа (където "X" са в горния ляв триъгълник) ще постави "X" в долния десен триъгълник, докато вашата мрежа има "X" в горния десен триъгълник.} \\
            }
        }
    }
    & & \\ \\

    \theutterance \stepcounter{utterance}  
    & & & \multicolumn{2}{p{0.3\linewidth}}{
        \cellcolor[rgb]{0.9,0.9,0.9}{
            \makecell[{{p{\linewidth}}}]{
                \texttt{\tiny{[GM$|$GM]}}
                \texttt{continue} \\
            }
        }
    }
    & & \\ \\

    \theutterance \stepcounter{utterance}  
    & & \multicolumn{4}{p{0.6\linewidth}}{
        \cellcolor[rgb]{0.9,0.9,0.9}{
            \makecell[{{p{\linewidth}}}]{
                \texttt{\tiny{[P1$\langle$GM]}}
                \texttt{Сега задайте въпрос, за да разберете нови аспекти на моята мрежа, които може да се различават от вашата мрежа. Започнете с "ВЪПРОС:" и не добавяйте нищо друго.} \\
            }
        }
    }
    & \\ \\

    \theutterance \stepcounter{utterance}  
    & \multicolumn{4}{p{0.6\linewidth}}{
        \cellcolor[rgb]{0.9,0.9,0.9}{
            \makecell[{{p{\linewidth}}}]{
                \texttt{\tiny{[P1$\rangle$GM]}}
                \texttt{ВЪПРОС: Ако отразите моята мрежа хоризонтално (огледално спрямо вертикалната ос), дали ще получите точно вашата мрежа?} \\
            }
        }
    }
    & & \\ \\

    \theutterance \stepcounter{utterance}  
    & & & \multicolumn{2}{p{0.3\linewidth}}{
        \cellcolor[rgb]{0.9,0.9,0.9}{
            \makecell[{{p{\linewidth}}}]{
                \texttt{\tiny{[GM$|$GM]}}
                \texttt{continue} \\
            }
        }
    }
    & & \\ \\

    \theutterance \stepcounter{utterance}  
    & & \multicolumn{4}{p{0.6\linewidth}}{
        \cellcolor[rgb]{0.9,0.9,0.9}{
            \makecell[{{p{\linewidth}}}]{
                \texttt{\tiny{[GM$\rangle$P2]}}
                \texttt{ОТГОВОР: Не, ако завъртите моята мрежа на 180 градуса, няма да получите вашата мрежа. Завъртането на 180 градуса на моята мрежа (където "X" са в горния ляв триъгълник) ще постави "X" в долния десен триъгълник, докато вашата мрежа има "X" в горния десен триъгълник.} \\
\texttt{ВЪПРОС: Ако отразите моята мрежа хоризонтално (огледално спрямо вертикалната ос), дали ще получите точно вашата мрежа?Започнете отговора си с "ОТГОВОР:" и не добавяйте нищо друго.} \\
            }
        }
    }
    & \\ \\

    \theutterance \stepcounter{utterance}  
    & & & \multicolumn{4}{p{0.6\linewidth}}{
        \cellcolor[rgb]{0.9,0.9,0.9}{
            \makecell[{{p{\linewidth}}}]{
                \texttt{\tiny{[GM$\langle$P2]}}
                \texttt{ОТГОВОР: Не, ако отразите моята мрежа хоризонтално, няма да получите моята мрежа. Хоризонталното отражение на вашата мрежа (с "X" в горния ляв триъгълник) ще постави "X" в горния десен триъгълник, което съвпада с моята мрежа, но въпросът беше дали ще получите моята мрежа чрез отражение на вашата. Тъй като вече установихме, че вашата мрежа е транспонираната версия на моята, а не хоризонтално отразената, отговорът е не.} \\
            }
        }
    }
     \\ \\

    \theutterance \stepcounter{utterance}  
    & & & \multicolumn{2}{p{0.3\linewidth}}{
        \cellcolor[rgb]{0.9,0.9,0.9}{
            \makecell[{{p{\linewidth}}}]{
                \texttt{\tiny{[GM$|$GM]}}
                \texttt{continue} \\
            }
        }
    }
    & & \\ \\

    \theutterance \stepcounter{utterance}  
    & & \multicolumn{4}{p{0.6\linewidth}}{
        \cellcolor[rgb]{0.9,0.9,0.9}{
            \makecell[{{p{\linewidth}}}]{
                \texttt{\tiny{[GM$\rangle$P2]}}
                \texttt{Сега задайте въпрос, за да разберете нови аспекти на моята мрежа, които може да се различават от вашата мрежа. Започнете с "ВЪПРОС:" и не добавяйте нищо друго.} \\
            }
        }
    }
    & \\ \\

    \theutterance \stepcounter{utterance}  
    & & & \multicolumn{4}{p{0.6\linewidth}}{
        \cellcolor[rgb]{0.9,0.9,0.9}{
            \makecell[{{p{\linewidth}}}]{
                \texttt{\tiny{[GM$\langle$P2]}}
                \texttt{ВЪПРОС: Ако отразите моята мрежа вертикално (огледално спрямо хоризонталната ос), дали ще получите точно моята мрежа?} \\
            }
        }
    }
     \\ \\

    \theutterance \stepcounter{utterance}  
    & & & \multicolumn{2}{p{0.3\linewidth}}{
        \cellcolor[rgb]{0.9,0.9,0.9}{
            \makecell[{{p{\linewidth}}}]{
                \texttt{\tiny{[GM$|$GM]}}
                \texttt{continue} \\
            }
        }
    }
    & & \\ \\

    \theutterance \stepcounter{utterance}  
    & & \multicolumn{4}{p{0.6\linewidth}}{
        \cellcolor[rgb]{0.9,0.9,0.9}{
            \makecell[{{p{\linewidth}}}]{
                \texttt{\tiny{[P1$\langle$GM]}}
                \texttt{ОТГОВОР: Не, ако отразите моята мрежа хоризонтално, няма да получите моята мрежа. Хоризонталното отражение на вашата мрежа (с "X" в горния ляв триъгълник) ще постави "X" в горния десен триъгълник, което съвпада с моята мрежа, но въпросът беше дали ще получите моята мрежа чрез отражение на вашата. Тъй като вече установихме, че вашата мрежа е транспонираната версия на моята, а не хоризонтално отразената, отговорът е не.} \\
\texttt{ВЪПРОС: Ако отразите моята мрежа вертикално (огледално спрямо хоризонталната ос), дали ще получите точно моята мрежа?Започнете отговора си с "ОТГОВОР:" и не добавяйте нищо друго.} \\
            }
        }
    }
    & \\ \\

    \theutterance \stepcounter{utterance}  
    & \multicolumn{4}{p{0.6\linewidth}}{
        \cellcolor[rgb]{0.9,0.9,0.9}{
            \makecell[{{p{\linewidth}}}]{
                \texttt{\tiny{[P1$\rangle$GM]}}
                \texttt{ОТГОВОР: Не, ако отразите моята мрежа вертикално, няма да получите моята мрежа. Вертикалното отражение на моята мрежа (с "X" в горния ляв триъгълник) ще постави "X" в долния ляв триъгълник, докато моята мрежа има "X" в горния ляв триъгълник.} \\
            }
        }
    }
    & & \\ \\

    \theutterance \stepcounter{utterance}  
    & & & \multicolumn{2}{p{0.3\linewidth}}{
        \cellcolor[rgb]{0.9,0.9,0.9}{
            \makecell[{{p{\linewidth}}}]{
                \texttt{\tiny{[GM$|$GM]}}
                \texttt{continue} \\
            }
        }
    }
    & & \\ \\

    \theutterance \stepcounter{utterance}  
    & & \multicolumn{4}{p{0.6\linewidth}}{
        \cellcolor[rgb]{0.9,0.9,0.9}{
            \makecell[{{p{\linewidth}}}]{
                \texttt{\tiny{[P1$\langle$GM]}}
                \texttt{Сега задайте въпрос, за да разберете нови аспекти на моята мрежа, които може да се различават от вашата мрежа. Започнете с "ВЪПРОС:" и не добавяйте нищо друго.} \\
            }
        }
    }
    & \\ \\

    \theutterance \stepcounter{utterance}  
    & \multicolumn{4}{p{0.6\linewidth}}{
        \cellcolor[rgb]{0.9,0.9,0.9}{
            \makecell[{{p{\linewidth}}}]{
                \texttt{\tiny{[P1$\rangle$GM]}}
                \texttt{ВЪПРОС: Ако отразите моята мрежа спрямо второстепенния диагонал (от горния десен ъгъл към долния ляв), дали ще получите точно вашата мрежа?} \\
            }
        }
    }
    & & \\ \\

    \theutterance \stepcounter{utterance}  
    & & & \multicolumn{2}{p{0.3\linewidth}}{
        \cellcolor[rgb]{0.9,0.9,0.9}{
            \makecell[{{p{\linewidth}}}]{
                \texttt{\tiny{[GM$|$GM]}}
                \texttt{continue} \\
            }
        }
    }
    & & \\ \\

    \theutterance \stepcounter{utterance}  
    & & \multicolumn{4}{p{0.6\linewidth}}{
        \cellcolor[rgb]{0.9,0.9,0.9}{
            \makecell[{{p{\linewidth}}}]{
                \texttt{\tiny{[GM$\rangle$P2]}}
                \texttt{ОТГОВОР: Не, ако отразите моята мрежа вертикално, няма да получите моята мрежа. Вертикалното отражение на моята мрежа (с "X" в горния ляв триъгълник) ще постави "X" в долния ляв триъгълник, докато моята мрежа има "X" в горния ляв триъгълник.} \\
\texttt{ВЪПРОС: Ако отразите моята мрежа спрямо второстепенния диагонал (от горния десен ъгъл към долния ляв), дали ще получите точно вашата мрежа?Започнете отговора си с "ОТГОВОР:" и не добавяйте нищо друго.} \\
            }
        }
    }
    & \\ \\

    \theutterance \stepcounter{utterance}  
    & & & \multicolumn{4}{p{0.6\linewidth}}{
        \cellcolor[rgb]{0.9,0.9,0.9}{
            \makecell[{{p{\linewidth}}}]{
                \texttt{\tiny{[GM$\langle$P2]}}
                \texttt{ОТГОВОР: Не, ако отразите моята мрежа спрямо второстепенния диагонал, няма да получите моята мрежа. Огледалното отражение спрямо второстепенния диагонал на мрежа с "X" в горния ляв триъгълник ще постави "X" в долния десен триъгълник, докато моята мрежа има "X" в горния десен триъгълник.} \\
            }
        }
    }
     \\ \\

    \theutterance \stepcounter{utterance}  
    & & & \multicolumn{2}{p{0.3\linewidth}}{
        \cellcolor[rgb]{0.9,0.9,0.9}{
            \makecell[{{p{\linewidth}}}]{
                \texttt{\tiny{[GM$|$GM]}}
                \texttt{continue} \\
            }
        }
    }
    & & \\ \\

    \theutterance \stepcounter{utterance}  
    & & \multicolumn{4}{p{0.6\linewidth}}{
        \cellcolor[rgb]{0.9,0.9,0.9}{
            \makecell[{{p{\linewidth}}}]{
                \texttt{\tiny{[GM$\rangle$P2]}}
                \texttt{Сега вземи решение. Какво мислиш: твоят квадрат и другият описан квадрат еднакви ли са?  Напиши "РЕШЕНИЕ: еднаква мрежа" ако смяташ, че са един и същ квадрат, или "РЕШЕНИЕ: различни мрежи" ако смяташ, че са различни квадрати. Не добавяй нищо друго.} \\
            }
        }
    }
    & \\ \\

    \theutterance \stepcounter{utterance}  
    & & & \multicolumn{4}{p{0.6\linewidth}}{
        \cellcolor[rgb]{0.9,0.9,0.9}{
            \makecell[{{p{\linewidth}}}]{
                \texttt{\tiny{[GM$\langle$P2]}}
                \texttt{РЕШЕНИЕ: различни мрежи} \\
            }
        }
    }
     \\ \\

    \theutterance \stepcounter{utterance}  
    & & & \multicolumn{2}{p{0.3\linewidth}}{
        \cellcolor[rgb]{0.9,0.9,0.9}{
            \makecell[{{p{\linewidth}}}]{
                \texttt{\tiny{[GM$|$GM]}}
                \texttt{continue} \\
            }
        }
    }
    & & \\ \\

    \theutterance \stepcounter{utterance}  
    & & & \multicolumn{2}{p{0.3\linewidth}}{
        \cellcolor[rgb]{0.9,0.9,0.9}{
            \makecell[{{p{\linewidth}}}]{
                \texttt{\tiny{[GM$|$GM]}}
                \texttt{success} \\
            }
        }
    }
    & & \\ \\

    \theutterance \stepcounter{utterance}  
    & & \multicolumn{4}{p{0.6\linewidth}}{
        \cellcolor[rgb]{0.9,0.9,0.9}{
            \makecell[{{p{\linewidth}}}]{
                \texttt{\tiny{[P1$\langle$GM]}}
                \texttt{ОТГОВОР: Не, ако отразите моята мрежа спрямо второстепенния диагонал, няма да получите моята мрежа. Огледалното отражение спрямо второстепенния диагонал на мрежа с "X" в горния ляв триъгълник ще постави "X" в долния десен триъгълник, докато моята мрежа има "X" в горния десен триъгълник.} \\
\texttt{Сега вземи решение. Какво мислиш: твоят квадрат и другият описан квадрат еднакви ли са?  Напиши "РЕШЕНИЕ: еднаква мрежа" ако смяташ, че са един и същ квадрат, или "РЕШЕНИЕ: различни мрежи" ако смяташ, че са различни квадрати. Не добавяй нищо друго.} \\
            }
        }
    }
    & \\ \\

    \theutterance \stepcounter{utterance}  
    & \multicolumn{4}{p{0.6\linewidth}}{
        \cellcolor[rgb]{0.9,0.9,0.9}{
            \makecell[{{p{\linewidth}}}]{
                \texttt{\tiny{[P1$\rangle$GM]}}
                \texttt{РЕШЕНИЕ: различни мрежи} \\
            }
        }
    }
    & & \\ \\

    \theutterance \stepcounter{utterance}  
    & & & \multicolumn{2}{p{0.3\linewidth}}{
        \cellcolor[rgb]{0.9,0.9,0.9}{
            \makecell[{{p{\linewidth}}}]{
                \texttt{\tiny{[GM$|$GM]}}
                \texttt{continue} \\
            }
        }
    }
    & & \\ \\

    \theutterance \stepcounter{utterance}  
    & & & \multicolumn{2}{p{0.3\linewidth}}{
        \cellcolor[rgb]{0.9,0.9,0.9}{
            \makecell[{{p{\linewidth}}}]{
                \texttt{\tiny{[GM$|$GM]}}
                \texttt{success} \\
            }
        }
    }
    & & \\ \\

\end{supertabular}
}

\end{document}
